% Options for packages loaded elsewhere
\PassOptionsToPackage{unicode}{hyperref}
\PassOptionsToPackage{hyphens}{url}
\documentclass[
  english,
  jou]{apa6}
\usepackage{xcolor}
\usepackage{amsmath,amssymb}
\setcounter{secnumdepth}{-\maxdimen} % remove section numbering
\usepackage{iftex}
\ifPDFTeX
  \usepackage[T1]{fontenc}
  \usepackage[utf8]{inputenc}
  \usepackage{textcomp} % provide euro and other symbols
\else % if luatex or xetex
  \usepackage{unicode-math} % this also loads fontspec
  \defaultfontfeatures{Scale=MatchLowercase}
  \defaultfontfeatures[\rmfamily]{Ligatures=TeX,Scale=1}
\fi
\usepackage{lmodern}
\ifPDFTeX\else
  % xetex/luatex font selection
\fi
% Use upquote if available, for straight quotes in verbatim environments
\IfFileExists{upquote.sty}{\usepackage{upquote}}{}
\IfFileExists{microtype.sty}{% use microtype if available
  \usepackage[]{microtype}
  \UseMicrotypeSet[protrusion]{basicmath} % disable protrusion for tt fonts
}{}
\makeatletter
\@ifundefined{KOMAClassName}{% if non-KOMA class
  \IfFileExists{parskip.sty}{%
    \usepackage{parskip}
  }{% else
    \setlength{\parindent}{0pt}
    \setlength{\parskip}{6pt plus 2pt minus 1pt}}
}{% if KOMA class
  \KOMAoptions{parskip=half}}
\makeatother
% Make \paragraph and \subparagraph free-standing
\makeatletter
\ifx\paragraph\undefined\else
  \let\oldparagraph\paragraph
  \renewcommand{\paragraph}{
    \@ifstar
      \xxxParagraphStar
      \xxxParagraphNoStar
  }
  \newcommand{\xxxParagraphStar}[1]{\oldparagraph*{#1}\mbox{}}
  \newcommand{\xxxParagraphNoStar}[1]{\oldparagraph{#1}\mbox{}}
\fi
\ifx\subparagraph\undefined\else
  \let\oldsubparagraph\subparagraph
  \renewcommand{\subparagraph}{
    \@ifstar
      \xxxSubParagraphStar
      \xxxSubParagraphNoStar
  }
  \newcommand{\xxxSubParagraphStar}[1]{\oldsubparagraph*{#1}\mbox{}}
  \newcommand{\xxxSubParagraphNoStar}[1]{\oldsubparagraph{#1}\mbox{}}
\fi
\makeatother
\usepackage{color}
\usepackage{fancyvrb}
\newcommand{\VerbBar}{|}
\newcommand{\VERB}{\Verb[commandchars=\\\{\}]}
\DefineVerbatimEnvironment{Highlighting}{Verbatim}{commandchars=\\\{\}}
% Add ',fontsize=\small' for more characters per line
\usepackage{framed}
\definecolor{shadecolor}{RGB}{248,248,248}
\newenvironment{Shaded}{\begin{snugshade}}{\end{snugshade}}
\newcommand{\AlertTok}[1]{\textcolor[rgb]{0.94,0.16,0.16}{#1}}
\newcommand{\AnnotationTok}[1]{\textcolor[rgb]{0.56,0.35,0.01}{\textbf{\textit{#1}}}}
\newcommand{\AttributeTok}[1]{\textcolor[rgb]{0.13,0.29,0.53}{#1}}
\newcommand{\BaseNTok}[1]{\textcolor[rgb]{0.00,0.00,0.81}{#1}}
\newcommand{\BuiltInTok}[1]{#1}
\newcommand{\CharTok}[1]{\textcolor[rgb]{0.31,0.60,0.02}{#1}}
\newcommand{\CommentTok}[1]{\textcolor[rgb]{0.56,0.35,0.01}{\textit{#1}}}
\newcommand{\CommentVarTok}[1]{\textcolor[rgb]{0.56,0.35,0.01}{\textbf{\textit{#1}}}}
\newcommand{\ConstantTok}[1]{\textcolor[rgb]{0.56,0.35,0.01}{#1}}
\newcommand{\ControlFlowTok}[1]{\textcolor[rgb]{0.13,0.29,0.53}{\textbf{#1}}}
\newcommand{\DataTypeTok}[1]{\textcolor[rgb]{0.13,0.29,0.53}{#1}}
\newcommand{\DecValTok}[1]{\textcolor[rgb]{0.00,0.00,0.81}{#1}}
\newcommand{\DocumentationTok}[1]{\textcolor[rgb]{0.56,0.35,0.01}{\textbf{\textit{#1}}}}
\newcommand{\ErrorTok}[1]{\textcolor[rgb]{0.64,0.00,0.00}{\textbf{#1}}}
\newcommand{\ExtensionTok}[1]{#1}
\newcommand{\FloatTok}[1]{\textcolor[rgb]{0.00,0.00,0.81}{#1}}
\newcommand{\FunctionTok}[1]{\textcolor[rgb]{0.13,0.29,0.53}{\textbf{#1}}}
\newcommand{\ImportTok}[1]{#1}
\newcommand{\InformationTok}[1]{\textcolor[rgb]{0.56,0.35,0.01}{\textbf{\textit{#1}}}}
\newcommand{\KeywordTok}[1]{\textcolor[rgb]{0.13,0.29,0.53}{\textbf{#1}}}
\newcommand{\NormalTok}[1]{#1}
\newcommand{\OperatorTok}[1]{\textcolor[rgb]{0.81,0.36,0.00}{\textbf{#1}}}
\newcommand{\OtherTok}[1]{\textcolor[rgb]{0.56,0.35,0.01}{#1}}
\newcommand{\PreprocessorTok}[1]{\textcolor[rgb]{0.56,0.35,0.01}{\textit{#1}}}
\newcommand{\RegionMarkerTok}[1]{#1}
\newcommand{\SpecialCharTok}[1]{\textcolor[rgb]{0.81,0.36,0.00}{\textbf{#1}}}
\newcommand{\SpecialStringTok}[1]{\textcolor[rgb]{0.31,0.60,0.02}{#1}}
\newcommand{\StringTok}[1]{\textcolor[rgb]{0.31,0.60,0.02}{#1}}
\newcommand{\VariableTok}[1]{\textcolor[rgb]{0.00,0.00,0.00}{#1}}
\newcommand{\VerbatimStringTok}[1]{\textcolor[rgb]{0.31,0.60,0.02}{#1}}
\newcommand{\WarningTok}[1]{\textcolor[rgb]{0.56,0.35,0.01}{\textbf{\textit{#1}}}}
\usepackage{graphicx}
\makeatletter
\newsavebox\pandoc@box
\newcommand*\pandocbounded[1]{% scales image to fit in text height/width
  \sbox\pandoc@box{#1}%
  \Gscale@div\@tempa{\textheight}{\dimexpr\ht\pandoc@box+\dp\pandoc@box\relax}%
  \Gscale@div\@tempb{\linewidth}{\wd\pandoc@box}%
  \ifdim\@tempb\p@<\@tempa\p@\let\@tempa\@tempb\fi% select the smaller of both
  \ifdim\@tempa\p@<\p@\scalebox{\@tempa}{\usebox\pandoc@box}%
  \else\usebox{\pandoc@box}%
  \fi%
}
% Set default figure placement to htbp
\def\fps@figure{htbp}
\makeatother
\ifLuaTeX
\usepackage[bidi=basic]{babel}
\else
\usepackage[bidi=default]{babel}
\fi
% get rid of language-specific shorthands (see #6817):
\let\LanguageShortHands\languageshorthands
\def\languageshorthands#1{}
\ifLuaTeX
  \usepackage[english]{selnolig} % disable illegal ligatures
\fi
\setlength{\emergencystretch}{3em} % prevent overfull lines
\providecommand{\tightlist}{%
  \setlength{\itemsep}{0pt}\setlength{\parskip}{0pt}}
% Manuscript styling
\usepackage{upgreek}
\captionsetup{font=singlespacing,justification=justified}

% Table formatting
\usepackage{longtable}
\usepackage{lscape}
% \usepackage[counterclockwise]{rotating}   % Landscape page setup for large tables
\usepackage{multirow}		% Table styling
\usepackage{tabularx}		% Control Column width
\usepackage[flushleft]{threeparttable}	% Allows for three part tables with a specified notes section
\usepackage{threeparttablex}            % Lets threeparttable work with longtable

% Create new environments so endfloat can handle them
% \newenvironment{ltable}
%   {\begin{landscape}\centering\begin{threeparttable}}
%   {\end{threeparttable}\end{landscape}}
\newenvironment{lltable}{\begin{landscape}\centering\begin{ThreePartTable}}{\end{ThreePartTable}\end{landscape}}

% Enables adjusting longtable caption width to table width
% Solution found at http://golatex.de/longtable-mit-caption-so-breit-wie-die-tabelle-t15767.html
\makeatletter
\newcommand\LastLTentrywidth{1em}
\newlength\longtablewidth
\setlength{\longtablewidth}{1in}
\newcommand{\getlongtablewidth}{\begingroup \ifcsname LT@\roman{LT@tables}\endcsname \global\longtablewidth=0pt \renewcommand{\LT@entry}[2]{\global\advance\longtablewidth by ##2\relax\gdef\LastLTentrywidth{##2}}\@nameuse{LT@\roman{LT@tables}} \fi \endgroup}

% \setlength{\parindent}{0.5in}
% \setlength{\parskip}{0pt plus 0pt minus 0pt}

% Overwrite redefinition of paragraph and subparagraph by the default LaTeX template
% See https://github.com/crsh/papaja/issues/292
\makeatletter
\renewcommand{\paragraph}{\@startsection{paragraph}{4}{\parindent}%
  {0\baselineskip \@plus 0.2ex \@minus 0.2ex}%
  {-1em}%
  {\normalfont\normalsize\bfseries\itshape\typesectitle}}

\renewcommand{\subparagraph}[1]{\@startsection{subparagraph}{5}{1em}%
  {0\baselineskip \@plus 0.2ex \@minus 0.2ex}%
  {-\z@\relax}%
  {\normalfont\normalsize\itshape\hspace{\parindent}{#1}\textit{\addperi}}{\relax}}
\makeatother

\makeatletter
\usepackage{etoolbox}
\patchcmd{\maketitle}
  {\section{\normalfont\normalsize\abstractname}}
  {\section*{\normalfont\normalsize\abstractname}}
  {}{\typeout{Failed to patch abstract.}}
\patchcmd{\maketitle}
  {\section{\protect\normalfont{\@title}}}
  {\section*{\protect\normalfont{\@title}}}
  {}{\typeout{Failed to patch title.}}
\makeatother

\usepackage{xpatch}
\makeatletter
\xapptocmd\appendix
  {\xapptocmd\section
    {\addcontentsline{toc}{section}{\appendixname\ifoneappendix\else~\theappendix\fi: #1}}
    {}{\InnerPatchFailed}%
  }
{}{\PatchFailed}
\makeatother
\keywords{self-prioritization, confirmatory study, reproduction\newline\indent Word count: }
\usepackage{dblfloatfix}


\usepackage{csquotes}
\usepackage{bookmark}
\IfFileExists{xurl.sty}{\usepackage{xurl}}{} % add URL line breaks if available
\urlstyle{same}
\hypersetup{
  pdftitle={Reproducing Hu et al.~(2020) Confirmatory Study},
  pdfauthor={Your Name1},
  pdflang={en-EN},
  pdfkeywords={self-prioritization, confirmatory study, reproduction},
  hidelinks,
  pdfcreator={LaTeX via pandoc}}

\title{Reproducing Hu et al.~(2020) Confirmatory Study}
\author{Your Name\textsuperscript{1}}
\date{2026-03-05}


\shorttitle{Good Me Bad Me (Confirmatory)}

\affiliation{\vspace{0.5cm}\textsuperscript{1} Your Institution}

\begin{document}
\maketitle

\section{Method}\label{method}

Describe participants, tasks, and procedure for the confirmatory study. Use the official article and supplementary material for exact details.

\section{Results}\label{results}

\begin{Shaded}
\begin{Highlighting}[]
\NormalTok{data\_dir }\OtherTok{\textless{}{-}} \FunctionTok{file.path}\NormalTok{(}\StringTok{".."}\NormalTok{, }\StringTok{"data"}\NormalTok{, }\StringTok{"processed"}\NormalTok{)}
\NormalTok{match\_wide }\OtherTok{\textless{}{-}} \FunctionTok{read.csv}\NormalTok{(}\FunctionTok{file.path}\NormalTok{(data\_dir, }\StringTok{"MS\_match\_behav\_wide.csv"}\NormalTok{), }\AttributeTok{stringsAsFactors =} \ConstantTok{FALSE}\NormalTok{)}
\NormalTok{categ\_wide }\OtherTok{\textless{}{-}} \FunctionTok{read.csv}\NormalTok{(}\FunctionTok{file.path}\NormalTok{(data\_dir, }\StringTok{"MS\_categ\_behav\_wide.csv"}\NormalTok{), }\AttributeTok{stringsAsFactors =} \ConstantTok{FALSE}\NormalTok{)}
\NormalTok{categ\_notask }\OtherTok{\textless{}{-}} \FunctionTok{read.csv}\NormalTok{(}\FunctionTok{file.path}\NormalTok{(data\_dir, }\StringTok{"MS\_categ\_behav\_noTask\_wide.csv"}\NormalTok{), }\AttributeTok{stringsAsFactors =} \ConstantTok{FALSE}\NormalTok{)}
\NormalTok{cross\_task }\OtherTok{\textless{}{-}} \FunctionTok{read.csv}\NormalTok{(}\FunctionTok{file.path}\NormalTok{(data\_dir, }\StringTok{"MS\_cross\_taskeffect\_wide.csv"}\NormalTok{), }\AttributeTok{stringsAsFactors =} \ConstantTok{FALSE}\NormalTok{)}
\NormalTok{jasp\_results }\OtherTok{\textless{}{-}} \FunctionTok{readRDS}\NormalTok{(}\FunctionTok{file.path}\NormalTok{(data\_dir, }\StringTok{"jasp\_reproduction"}\NormalTok{, }\StringTok{"jasp\_reproduction\_results.rds"}\NormalTok{))}
\end{Highlighting}
\end{Shaded}

Summarize the confirmatory study analyses here, mirroring reported results.

\begin{Shaded}
\begin{Highlighting}[]
\NormalTok{compare\_dir }\OtherTok{\textless{}{-}} \FunctionTok{file.path}\NormalTok{(data\_dir, }\StringTok{"jasp\_reproduction"}\NormalTok{)}
\NormalTok{compare\_ttests }\OtherTok{\textless{}{-}} \FunctionTok{file.path}\NormalTok{(compare\_dir, }\StringTok{"compare\_ttests.csv"}\NormalTok{)}
\NormalTok{compare\_anova }\OtherTok{\textless{}{-}} \FunctionTok{file.path}\NormalTok{(compare\_dir, }\StringTok{"compare\_anova.csv"}\NormalTok{)}
\NormalTok{compare\_cor }\OtherTok{\textless{}{-}} \FunctionTok{file.path}\NormalTok{(compare\_dir, }\StringTok{"compare\_cor.csv"}\NormalTok{)}

\NormalTok{safe\_read\_csv }\OtherTok{\textless{}{-}} \ControlFlowTok{function}\NormalTok{(path) \{}
  \ControlFlowTok{if}\NormalTok{ (}\SpecialCharTok{!}\FunctionTok{file.exists}\NormalTok{(path) }\SpecialCharTok{||} \FunctionTok{file.info}\NormalTok{(path)}\SpecialCharTok{$}\NormalTok{size }\SpecialCharTok{\textless{}=} \DecValTok{3}\NormalTok{) \{}
    \FunctionTok{return}\NormalTok{(}\ConstantTok{NULL}\NormalTok{)}
\NormalTok{  \}}
\NormalTok{  out }\OtherTok{\textless{}{-}} \FunctionTok{tryCatch}\NormalTok{(}\FunctionTok{read.csv}\NormalTok{(path, }\AttributeTok{stringsAsFactors =} \ConstantTok{FALSE}\NormalTok{), }\AttributeTok{error =} \ControlFlowTok{function}\NormalTok{(e) }\ConstantTok{NULL}\NormalTok{)}
  \ControlFlowTok{if}\NormalTok{ (}\FunctionTok{is.null}\NormalTok{(out) }\SpecialCharTok{||} \FunctionTok{nrow}\NormalTok{(out) }\SpecialCharTok{==} \DecValTok{0}\NormalTok{) \{}
    \FunctionTok{return}\NormalTok{(}\ConstantTok{NULL}\NormalTok{)}
\NormalTok{  \}}
\NormalTok{  out}
\NormalTok{\}}

\NormalTok{anova\_df }\OtherTok{\textless{}{-}} \FunctionTok{safe\_read\_csv}\NormalTok{(compare\_anova)}
\NormalTok{ttest\_df }\OtherTok{\textless{}{-}} \FunctionTok{safe\_read\_csv}\NormalTok{(compare\_ttests)}
\NormalTok{cor\_df }\OtherTok{\textless{}{-}} \FunctionTok{safe\_read\_csv}\NormalTok{(compare\_cor)}

\ControlFlowTok{if}\NormalTok{ (}\SpecialCharTok{!}\FunctionTok{is.null}\NormalTok{(anova\_df)) \{}
  \FunctionTok{cat}\NormalTok{(}\StringTok{"\#\# JASP comparison (ANOVA)}\SpecialCharTok{\textbackslash{}n}\StringTok{"}\NormalTok{)}
  \FunctionTok{print}\NormalTok{(utils}\SpecialCharTok{::}\FunctionTok{head}\NormalTok{(anova\_df))}
\NormalTok{\}}
\ControlFlowTok{if}\NormalTok{ (}\SpecialCharTok{!}\FunctionTok{is.null}\NormalTok{(ttest\_df)) \{}
  \FunctionTok{cat}\NormalTok{(}\StringTok{"\#\# JASP comparison (T{-}tests)}\SpecialCharTok{\textbackslash{}n}\StringTok{"}\NormalTok{)}
  \FunctionTok{print}\NormalTok{(utils}\SpecialCharTok{::}\FunctionTok{head}\NormalTok{(ttest\_df))}
\NormalTok{\}}
\ControlFlowTok{if}\NormalTok{ (}\SpecialCharTok{!}\FunctionTok{is.null}\NormalTok{(cor\_df)) \{}
  \FunctionTok{cat}\NormalTok{(}\StringTok{"\#\# JASP comparison (Correlations)}\SpecialCharTok{\textbackslash{}n}\StringTok{"}\NormalTok{)}
  \FunctionTok{print}\NormalTok{(utils}\SpecialCharTok{::}\FunctionTok{head}\NormalTok{(cor\_df))}
\NormalTok{\}}
\end{Highlighting}
\end{Shaded}

\subsection{JASP comparison (Correlations)}\label{jasp-comparison-correlations}

\begin{verbatim}
                    analysis                    var1
\end{verbatim}

1 MS\_cross\_taskeffect\_wide:0:cor d\_goodslf\_goodoth
2 MS\_cross\_taskeffect\_wide:0:cor d\_goodslf\_goodoth
3 MS\_cross\_taskeffect\_wide:0:cor Val\_ACC\_goodslf\_goodoth
4 MS\_cross\_taskeffect\_wide:2:cor d\_goodslf\_badslf
5 MS\_cross\_taskeffect\_wide:2:cor d\_goodslf\_badslf
6 MS\_cross\_taskeffect\_wide:2:cor Val\_ACC\_goodslf\_badslf
var2 r\_jasp r\_r r\_diff p\_jasp
1 Val\_ACC\_goodslf\_goodoth 0.20514179 0.20514179 -8.326673e-17 0.1982124
2 Id\_ACC\_goodslf\_goodoth 0.13063634 0.13063634 4.163336e-16 0.4155821
3 Id\_ACC\_goodslf\_goodoth 0.14413807 0.14413807 -2.220446e-16 0.3685996
4 Val\_ACC\_goodslf\_badslf 0.24164200 0.24164200 -3.885781e-16 0.1279914
5 Id\_ACC\_goodslf\_badslf 0.06007816 0.06007816 -4.163336e-17 0.7090527
6 Id\_ACC\_goodslf\_badslf 0.18571572 0.18571572 0.000000e+00 0.2450190
p\_r p\_diff
1 0.1982124 -2.775558e-17
2 0.4155821 5.551115e-16
3 0.3685996 -1.665335e-16
4 0.1279914 5.273559e-16
5 0.7090527 -2.220446e-16
6 0.2450190 -1.387779e-16

\begin{Shaded}
\begin{Highlighting}[]
\ControlFlowTok{if}\NormalTok{ (}\SpecialCharTok{!}\FunctionTok{is.null}\NormalTok{(jasp\_results}\SpecialCharTok{$}\NormalTok{r}\SpecialCharTok{$}\NormalTok{MS\_rep\_match\_behav)) \{}
  \FunctionTok{cat}\NormalTok{(}\StringTok{"\#\# Matching task (JASP{-}matched outputs)}\SpecialCharTok{\textbackslash{}n}\StringTok{"}\NormalTok{)}
  \FunctionTok{print}\NormalTok{(jasp\_results}\SpecialCharTok{$}\NormalTok{r}\SpecialCharTok{$}\NormalTok{MS\_rep\_match\_behav[[}\DecValTok{1}\NormalTok{]])}
\NormalTok{\}}
\end{Highlighting}
\end{Shaded}

\subsection{Matching task (JASP-matched outputs)}\label{matching-task-jasp-matched-outputs}

\subsection{Bayes factor analysis}\label{bayes-factor-analysis}

{[}1{]} MValence + Subject : 8308161 ±0.77\%
{[}2{]} ID + Subject : 0.5847399 ±4.84\%
{[}3{]} MValence + ID + Subject : 7105599 ±1.7\%
{[}4{]} MValence + ID + MValence:ID + Subject : 514807069 ±2.72\%

Against denominator:
value \textasciitilde{} Subject
---
Bayes factor type: BFlinearModel, JZS

\section{Discussion}\label{discussion}

Summarize the reproduction outcome and note any deviations from the original pipeline.


\end{document}
